% =====================================================================
%  06_referencia.tex — Variables, API, convenciones y base técnica
% =====================================================================
\chapter{Referencia}
\label{cap:referencia}

\section{Variables de entorno}
\label{sec:variables}

Esta tabla resume el \path{.env}. La fuente de verdad siempre es
\path{.env.example}: si agregas una variable, actualiza los tres lugares
(\path{config.py}, \path{.env.example} y esta tabla).

{\footnotesize
\begin{longtable}{@{}L{4.3cm}L{3.4cm}L{6.3cm}@{}}
  \toprule
  \textbf{Variable} & \textbf{Default} & \textbf{Qué hace} \\
  \midrule
  \endhead
  \multicolumn{3}{@{}l}{\textbf{\color{azulCISA}Cámara}} \\
  \addlinespace[2pt]
  \texttt{RTSP\_URL} & --- & Fuente activa: cámara o \texttt{.mp4} de prueba. \\
  \texttt{RTSP\_URL\_CAM2} & vacío & Segunda cámara opcional (botón ``Cam 2''). \\
  \texttt{CAM\_BACKEND} & vacío & \texttt{ffmpeg} desactiva la decodificación por
    hardware. \\
  \texttt{VIDEOS\_PRUEBA} & vacío & Carpetas o archivos de prueba, separados por
    coma. \\
  \midrule
  \multicolumn{3}{@{}l}{\textbf{\color{azulCISA}Modelos}} \\
  \addlinespace[2pt]
  \texttt{MODEL\_PATH} & --- & Motor M2 (zona del número). \\
  \texttt{DIGITS\_MODEL\_PATH} & --- & Motor M3 (dígitos). \\
  \texttt{AUTOBUS\_MODEL\_PATH} & --- & Motor M1 (autobús). \\
  \midrule
  \multicolumn{3}{@{}l}{\textbf{\color{azulCISA}Rendimiento}} \\
  \addlinespace[2pt]
  \texttt{YOLO\_FPS\_MAX} & 0 & Cuadros por segundo de M1; 0 = sin límite. \\
  \texttt{M2\_FPS\_MAX} & 5 & Cuadros por segundo de la lectura; bajarlo ahorra
    GPU. \\
  \texttt{CROPS\_CADA\_SEG} & 1 & Guarda un crop cada $N$ segundos por camión
    (0 = todos). \\
  \midrule
  \multicolumn{3}{@{}l}{\textbf{\color{azulCISA}Línea virtual y registro}} \\
  \addlinespace[2pt]
  \texttt{LINEA\_P1}, \texttt{LINEA\_P2} & --- & Puntos de la línea virtual. \\
  \texttt{LADO\_ENTRADA} & izquierda & Lado que cuenta como ``adentro''. \\
  \texttt{MARGEN\_LINEA\_PX} & 30 & Franja donde el lado no cambia. \\
  \texttt{MANIOBRA\_SEG} & 20 & Regreso antes de este tiempo = maniobra. \\
  \texttt{DEBOUNCE\_SEG} & 60 & Tiempo mínimo entre dos registros del mismo
    número. \\
  \texttt{ESPERA\_FIN\_SEG} & 2 & Segundos sin ver al camión para darlo por
    pasado. \\
  \texttt{ESPERA\_MAX\_SEG} & 25 & Máximo a esperar para registrar; debe ser
    mayor que \texttt{MANIOBRA\_SEG}. \\
  \texttt{VOTOS\_FUERA\_PADRON} & 4 & Lecturas iguales para registrar un número
    que no está en el padrón. \\
  \texttt{FOTOS\_POR\_AVISO} & 3 & Fotos por aviso (las mejores, de momentos
    distintos). \\
  \texttt{FOTOS\_SEPARACION\_SEG} & 1.0 & Segundos mínimos entre foto y foto. \\
  \texttt{ID\_PUERTA} & --- & Puerta del patio (1 o 2). \\
  \midrule
  \multicolumn{3}{@{}l}{\textbf{\color{azulCISA}Avisos y validación}} \\
  \addlinespace[2pt]
  \texttt{TG\_TOKEN} & --- & Token del bot de Telegram. \\
  \texttt{TG\_CHAT\_ID} & --- & Grupo que recibe los avisos. \\
  \texttt{TG\_VALIDADOR\_IDS} & --- & Usuarios que pueden validar por Telegram. \\
  \texttt{SIN\_TELEGRAM} & 0 & \texttt{1} = no manda nada (pruebas). \\
  \texttt{GEMINI\_API\_KEY} & vacío & Clave de Gemini (solo en el \path{.env}). \\
  \texttt{GEMINI\_MODEL} & flash-lite & Modelo de Gemini que revisa la foto. \\
  \texttt{VALIDAR\_CON\_GEMINI} & sigue a Telegram & \texttt{1} = revisar aunque
    Telegram esté apagado. \\
  \texttt{WHATSAPP\_NUMERO} & vacío & Número que recibe los avisos (52 + 10
    dígitos); vacío = apagado. \\
  \texttt{WHATSAPP\_EMISOR} & vacío & Teléfono que manda; si está, vincula con
    código de 8 letras en vez de QR. \\
  \texttt{WHATSAPP\_URL} & \url{http://127.0.0.1:5002} & Puerto del servicio de
    WhatsApp. \\
  \midrule
  \multicolumn{3}{@{}l}{\textbf{\color{azulCISA}Carpetas}} \\
  \addlinespace[2pt]
  \texttt{DB\_PATH} & \path{datos/detecciones.db} & Base de datos SQLite. \\
  \texttt{FRAMES\_DIR} & --- & Fotos de los avisos. \\
  \texttt{CROPS\_DIR} & \path{/media/cisa/JETSON_SD/crops_dataset} & Recortes
    para revisión y dataset. \\
  \texttt{GRABACIONES\_DIR} & \path{/media/cisa/JETSON_SD} & Base de
    \path{cam1_grabacion}. \\
  \texttt{PADRON\_PATH} & \texttt{padron.csv} & Padrón de unidades. \\
  \midrule
  \multicolumn{3}{@{}l}{\textbf{\color{azulCISA}Web}} \\
  \addlinespace[2pt]
  \texttt{FLASK\_HOST} & --- & Interfaz donde escucha Flask. \\
  \texttt{FLASK\_PORT} & --- & Puerto (5001 en la operación normal). \\
  \texttt{MJPEG\_FPS} & 10 & Cuadros por segundo del video en vivo. \\
  \texttt{MJPEG\_ANCHO} & 960 & Ancho en píxeles del video en vivo. \\
  \texttt{LISTA\_DETECCIONES} & 10 & Detecciones por página en el panel. \\
  \bottomrule
  \caption{Variables de entorno del sistema.}
  \label{tab:env}
\end{longtable}
}

\section{API HTTP}

Las rutas viven todas en \path{web/rutas.py}. Devuelven JSON salvo las que
sirven páginas, video o archivos.

{\small
\begin{longtable}{@{}L{5.4cm}L{8.6cm}@{}}
  \toprule
  \textbf{Ruta} & \textbf{Devuelve} \\
  \midrule
  \endhead
  \texttt{GET /} & Tablero principal. \\
  \texttt{GET /livevideo} & Página de video en vivo y detecciones. \\
  \texttt{GET /grabaciones} & Página de grabaciones. \\
  \texttt{GET /home} & Página de inicio alternativa. \\
  \texttt{GET /health} & Estado del pipeline: conexión, FPS, cámara activa. \\
  \texttt{GET /api/estado} & Estado completo para el tablero. \\
  \texttt{GET /api/dashboard} & Datos agregados del tablero. \\
  \texttt{GET /api/detecciones} & Detecciones con búsqueda, empresa, fechas y
    paginación. \\
  \texttt{GET /api/detecciones-recent} & Últimas detecciones. \\
  \texttt{GET /api/detecciones/empresas} & Empresas disponibles para el filtro. \\
  \texttt{GET /api/estancias} & Estancias e historial. \\
  \texttt{GET /api/snapshot} & Foto del último cuadro. \\
  \texttt{GET /api/video.mjpeg} & Video en vivo por MJPEG. \\
  \texttt{POST /api/switch-camera} & Cambia entre cámara y cámara 2. \\
  \texttt{GET /api/grabaciones} & Lista de grabaciones por cámara. \\
  \texttt{GET /frames/<archivo>} & Foto de un aviso. \\
  \texttt{GET /thumb/<archivo>} & Miniatura de una foto. \\
  \path{GET /grabaciones/file/<cam>/<archivo>} & Archivo de video. \\
  \texttt{GET /api/crops/stats} & Estadísticas de recortes guardados. \\
  \bottomrule
  \caption{Rutas HTTP del sistema.}
  \label{tab:api}
\end{longtable}
}

\section{Convenciones del código}

\begin{itemize}
  \item \textbf{La configuración se lee una sola vez} en \path{config.py}.
        Ningún otro módulo llama a \texttt{os.getenv}.
  \item \textbf{Las rutas web no tienen lógica de negocio}: leen de
        \path{logica/estado.py} y de \path{metricas/}, y responden.
  \item \textbf{Un esquema, una migración}: los cambios de base se hacen en
        \path{metricas/schema.sql} y se migran en \path{storage.py}.
  \item \textbf{Los umbrales de los modelos} son constantes de \path{config.py};
        lo que cambia por patio va en el \path{.env}.
  \item \textbf{Los comentarios explican el porqué}, no el qué.
\end{itemize}

\section{Cómo extender el sistema}

\begin{longtable}{@{}L{5cm}L{9cm}@{}}
  \toprule
  \textbf{Quiero\ldots} & \textbf{Tocar} \\
  \midrule
  \endhead
  Agregar una variable de configuración & \path{config.py},
    \path{.env.example} y la tabla~\ref{tab:env}. \\
  Agregar una ruta web & \path{web/rutas.py}; si necesita datos nuevos,
    \path{metricas/}. \\
  Cambiar la base de datos & \path{metricas/schema.sql} y la migración en
    \path{storage.py}. \\
  Agregar un aviso & \path{avisos/}; respetar \texttt{SIN\_TELEGRAM}. \\
  Cambiar la lógica de cruce & \path{ml/linea_virtual.py} y sus variables en
    \path{.env}. \\
  Cambiar la interfaz & \path{templates/} y \path{static/}. \\
  \bottomrule
  \caption{Guía rápida para cambios.}
  \label{tab:extender}
\end{longtable}

\section{Base técnica}

El sistema se apoya en herramientas abiertas y en el SDK de NVIDIA:

\begin{itemize}
  \item La detección usa YOLO~\cite{redmon2016yolo}, implementado con
        Ultralytics~\cite{ultralytics}.
  \item El seguimiento entre cuadros es ByteTrack~\cite{zhang2022bytetrack}.
  \item Los modelos corren en TensorRT sobre la Jetson~\cite{tensorrt,jetson}.
  \item La lectura de video usa OpenCV~\cite{opencv}.
  \item El servidor web es Flask~\cite{flask}, y la base es
        SQLite~\cite{sqlite}.
  \item Los avisos usan la API de bots de Telegram~\cite{telegram},
        la API de Gemini~\cite{gemini} y Baileys~\cite{baileys}.
\end{itemize}
